\section{Introduction}
\label{Sec:Introduction}

Termination analysis of infinite-state systems is a fundamental problem 
in formal methods.
Proving termination and nontermination over LIA is generally undecidable~\cite{DBLP:conf/cav/Braverman06}, 
yet many practical techniques succeed on large classes of real-world 
programs.

Modern techniques treat termination and nontermination as separate 
problems, since each has a different proof procedure.
State-of-the-art termination techniques centre around the synthesis of 
ranking functions for the transition relation~\cite{DBLP:conf/vmcai/PodelskiR04,DBLP:conf/cav/KuraUH20,DBLP:journals/jacm/Ben-AmramG14,DBLP:conf/cav/HeizmannHP14,DBLP:journals/jar/GieslABEFFHOPSS17,DBLP:conf/esop/SaritaSGSV26,DBLP:conf/sas/Ben-AmramDG19,DBLP:conf/tacas/UrbanGK16}.
The existence of a ranking function witnesses termination by providing a measure that strictly decreases along every 
execution.
Alternative approaches are based on the fixpoint computation~\cite{DBLP:journals/pacmpl/UnnoTGK23}, and the synthesis of disjunctively well-founded 
transition invariants~\cite{Transition-Invariants-Revisited-CAV26,DBLP:conf/pldi/CookPR06,DBLP:conf/cav/KroeningSTW10,DBLP:conf/tacas/TsitovichSWK11}.
Nontermination analysis techniques are mostly centered around the detection of recurrent 
sets — sets of states from which execution cannot escape~\cite{DBLP:conf/sas/Ben-AmramDG19,DBLP:conf/tacas/ChenCFNO14,DBLP:conf/cade/FrohnG22,DBLP:conf/tacas/BrockschmidtCIK16,DBLP:conf/foveoos/BrockschmidtSOG11,DBLP:conf/cav/LarrazNORR14,DBLP:journals/entcs/BiereAS02}.
There also exist alternative nontermination approaches based on Hoare-style reasoning~\cite{DBLP:conf/pldi/LeQC15}, safety proving~\cite{DBLP:conf/tacas/ChenCFNO14,DBLP:journals/entcs/BiereAS02}, or representing infinite runs as sums of geometric series~\cite{DBLP:conf/tacas/LeikeH18}.


In this work, we present a unified framework that is capable of proving both termination and nontermination.
This technique is constructed by extending the safety-based nontermination 
analysis~\cite{DBLP:conf/tacas/ChenCFNO14} (SNA), allowing it to additionally prove termination.
Implemented termination analysis is based on a new interpolation-based method for generation of
disjunctively well-founded transition invariants, which utilises the intermediate results
of nontermination analysis.
Proposed approach uses terminating traces constructed by SNA, overapproximating them with Craig's interpolation.
These overapproximations are used as transition invariant candidate, which can prove termination if 
it is a valid transition invariant over all reachable states.
Additionally, algorithm computes states over which candidate transition invariant is not valid.
Those states are separately analyzed for termination and nontermination, introducing smaller 
subproblem, compared to the analysis of the whole transition system.


Our new technique was implemented in Golem CHC Solver. 
We evaluated our technique over the Integer Transition Systems benchmarks from the Termination Competition. 
Proposed approach that combines termination and nontermination frameworks demonstrated significant 
improvement for nontermination analysis, compared to the original SNA algorithm, 
while also solving a large number of terminating instances.
The comparison with winners of Termination Competition 2025 (KoAT, LoAT~\cite{10.1007/978-3-031-90660-2_13}, T2~\cite{DBLP:conf/tacas/BrockschmidtCIK16}) 
demonstrated that the developed approach is competitive with the state-of-the-art tools, and is able to solve unique instances, including instances
that were never solved in the history of Termination competition.

Overall, this paper presents following contributions: (1) a new interpolation-based method for the synthesis of disjunctively well-founded transition invariants, 
(2) a unified termination/nontermination procedure, which utilises the intermediate results of analysis for self-guidance and (3) an implementation of the described 
procedure in Golem and evaluation of proposed approach.